\section{Konjunktur}

\subsection{Konjunktur vs. Wachstum}

\textbf{Konjunktur} beschreibt die kurzfristigen Schwankungen der wirtschaftlichen Aktivität 
um den langfristigen Wachstumspfad einer Volkswirtschaft.

\textbf{Wirtschaftswachstum} beschreibt dagegen die langfristige Zunahme des 
Produktionspotentials einer Volkswirtschaft.

Konjunktur wird typischerweise anhand von Veränderungen des realen BIP gemessen.

\textbf{Rezession}

\[
\text{Rezession} = \text{Rückgang des realen BIP während mindestens zwei Quartalen}
\]

Konjunkturschwankungen entstehen durch Veränderungen von:

\begin{itemize}
\item aggregierter Nachfrage (AN)
\item aggregiertem Angebot (AAK)
\end{itemize}


\subsection{Konjunkturindikatoren}

Konjunkturindikatoren liefern Informationen über die aktuelle oder zukünftige wirtschaftliche Entwicklung.

Beispiele:

\begin{itemize}
\item Preisentwicklung
\item Bestellungseingänge
\item Investitionsverhalten
\item Lohnentwicklung
\item Arbeitslosigkeit
\item Geldmenge
\item Wechselkurse
\item Konsumverhalten
\item Konsumentenstimmung
\item Export- und Importentwicklung
\item Zinsen
\end{itemize}

\textbf{Arten von Indikatoren}

\begin{itemize}
\item \textbf{Vorlaufende Indikatoren}  
zeigen zukünftige Entwicklungen frühzeitig (z.B. Konjunkturbarometer)

\item \textbf{Gleichlaufende Indikatoren}  
verändern sich gleichzeitig mit der Wirtschaft (z.B. Investitionen)

\item \textbf{Nachlaufende Indikatoren}  
reagieren erst nach der Konjunkturentwicklung (z.B. Löhne, Arbeitslosigkeit)
\end{itemize}


\subsection{Makroökonomisches Grundmodell}

Das makroökonomische Gleichgewicht entsteht im Schnittpunkt von:

\begin{itemize}
\item aggregierter Nachfrage (AN)
\item aggregiertem Angebot (AAK)
\end{itemize}

Das Gleichgewicht bestimmt:

\begin{itemize}
\item Preisniveau
\item reales BIP
\end{itemize}

Die \textbf{Kapazitätsgrenze} beschreibt die Produktionskapazität bei Normalbeschäftigung.


\subsection{Nachfrageschocks}

Nachfrageschocks verschieben die Kurve der aggregierten Nachfrage (AN).

\textbf{Negativer Nachfrageschock}

AN verschiebt sich nach links.

Mögliche Ursachen:

\begin{itemize}
\item Rückgang des Konsums
\item Rückgang der Investitionen
\item geringere Staatsausgaben
\item sinkende Exporte
\end{itemize}

Auswirkungen:

\begin{itemize}
\item reales BIP sinkt
\item Preisniveau sinkt
\item konjunkturelle Arbeitslosigkeit entsteht
\end{itemize}

\textbf{Positiver Nachfrageschock}

AN verschiebt sich nach rechts.

Mögliche Ursachen:

\begin{itemize}
\item steigender Konsum
\item höhere Investitionen
\item höhere Staatsausgaben
\item steigende Exporte
\end{itemize}

Auswirkungen:

\begin{itemize}
\item reales BIP steigt
\item Preisniveau steigt
\end{itemize}


\subsection{Angebotsschocks}

Angebotsschocks verändern die kurzfristige Angebotskurve (AAK).

\textbf{Negativer Angebotsschock}

AAK verschiebt sich nach links.

Mögliche Ursachen:

\begin{itemize}
\item steigende Rohstoffpreise
\item steigende Löhne
\item höhere Steuern
\item Produktionsstörungen
\end{itemize}

Auswirkungen:

\begin{itemize}
\item Preisniveau steigt
\item reales BIP sinkt
\end{itemize}

\textbf{Positiver Angebotsschock}

AAK verschiebt sich nach rechts.

Mögliche Ursachen:

\begin{itemize}
\item technologische Fortschritte
\item sinkende Produktionskosten
\item tiefere Steuern
\item effizientere Produktionsprozesse
\end{itemize}

Auswirkungen:

\begin{itemize}
\item reales BIP steigt
\item Preisniveau sinkt
\end{itemize}


\subsection{Konjunkturzyklen}

Konjunkturschwankungen können in vier Situationen eingeteilt werden:

\begin{itemize}
\item \textbf{Inflationärer Boom}  
reales BIP steigt, Preisniveau steigt

\item \textbf{Deflationärer Boom}  
reales BIP steigt, Preisniveau sinkt

\item \textbf{Stagflation}  
reales BIP sinkt, Preisniveau steigt

\item \textbf{Depression}  
reales BIP sinkt, Preisniveau sinkt
\end{itemize}


\subsection{Konjunkturelle Arbeitslosigkeit}

Konjunkturelle Arbeitslosigkeit entsteht bei einer Unterauslastung der Wirtschaft.

\[
\text{Arbeitssuchende} > \text{offene Stellen}
\]

Unterschied zu Sockelarbeitslosigkeit:

\begin{itemize}
\item \textbf{strukturelle Arbeitslosigkeit}  
Mismatch zwischen Qualifikation und Stellen

\item \textbf{friktionelle Arbeitslosigkeit}  
Sucharbeitslosigkeit beim Jobwechsel
\end{itemize}


\subsection{Konjunkturpolitik}

Ziel der Konjunkturpolitik ist die Stabilisierung der Wirtschaft.

Es existieren drei grundlegende Ansätze.

\textbf{1. Keine aktive Politik („Nichts tun“)}

Marktmechanismen führen langfristig zurück zum Gleichgewicht:

\begin{itemize}
\item sinkende Löhne
\item sinkende Produktionskosten
\item Verschiebung der Angebotskurve nach rechts
\end{itemize}

\textbf{2. Keynesianische Konjunkturpolitik}

Der Staat stimuliert die Nachfrage aktiv.

Instrumente:

\begin{itemize}
\item Fiskalpolitik (Staatsausgaben, Steuern)
\item Geldpolitik (Zinsen, Geldmenge)
\end{itemize}

Ziel:

\begin{itemize}
\item Konsum erhöhen
\item Investitionen erhöhen
\item Staatsausgaben erhöhen
\item Nettoexporte erhöhen
\end{itemize}

Problem:

\begin{itemize}
\item Wirkungsverzögerungen (time lags)
\end{itemize}


\subsection{Automatische Stabilisatoren}

Automatische Stabilisatoren wirken ohne aktive politische Entscheidung.

Beispiele:

\begin{itemize}
\item progressives Steuersystem
\item Arbeitslosenversicherung
\item staatliche Transfers
\end{itemize}

Diese stabilisieren die Nachfrage während Rezessionen automatisch.


\subsection{Schweizer Konjunkturpolitik}

Die Bundesverfassung verpflichtet den Staat zur Stabilisierung der Wirtschaft.

Ziele:

\begin{itemize}
\item Bekämpfung von Arbeitslosigkeit
\item Preisstabilität
\end{itemize}

Die Schweiz setzt vor allem auf automatische Stabilisatoren:

\begin{itemize}
\item Schuldenbremse
\item Arbeitslosenversicherung
\item progressives Steuersystem
\end{itemize}


\subsection{Konjunkturprogramm 2008--2010}

Als Reaktion auf die Finanzkrise beschloss der Bund mehrere 
Stabilisierungsprogramme.

Gesamtvolumen:

\[
\approx 2.1 \text{ Milliarden CHF}
\]

Massnahmen:

\begin{itemize}
\item Infrastrukturinvestitionen
\item Förderung erneuerbarer Energien
\item Bekämpfung der Arbeitslosigkeit
\item CO$_2$-Rückverteilung
\end{itemize}

Ziel der Programme:

\begin{itemize}
\item \textbf{timely} (rechtzeitig)
\item \textbf{targeted} (zielgerichtet)
\item \textbf{temporary} (vorübergehend)
\end{itemize}


\subsection{Corona-Krise 2020}

Zur Stabilisierung der Wirtschaft beschloss die Schweiz ein 
Massnahmenpaket von rund

\[
\approx 60 \text{ Milliarden CHF}
\]

Massnahmen:

\begin{itemize}
\item Liquiditätshilfen für Unternehmen
\item COVID-Überbrückungskredite
\item Kurzarbeitsentschädigung
\item Erwerbsausfallentschädigung
\item Unterstützung für Kultur- und Sportsektor
\end{itemize}


\subsection{Zyklische Wirtschaftskrisen}

Konjunkturzyklen sind ein wiederkehrendes Phänomen.

Grundprinzip:

\[
\text{Nach der Krise ist vor der Krise}
\]

Die Wirtschaft bewegt sich langfristig um ihren Wachstumspfad, 
wird jedoch kurzfristig durch Nachfrage- und Angebotsschocks beeinflusst.